
\documentclass{article}
\usepackage{colortbl}
\usepackage{makecell}
\usepackage{multirow}
\usepackage{supertabular}

\begin{document}

\newcounter{utterance}

\centering \large Interaction Transcript for game `cladder', experiment `full\_v1.5\_default', episode 9887 with Qwen3.6{-}27B.
\vspace{24pt}

{ \footnotesize  \setcounter{utterance}{1}
\setlength{\tabcolsep}{0pt}
\begin{supertabular}{c@{$\;$}|p{.15\linewidth}@{}p{.15\linewidth}p{.15\linewidth}p{.15\linewidth}p{.15\linewidth}p{.15\linewidth}}
    \# & \multicolumn{2}{c}{Player} && \multicolumn{2}{c}{Game Master} \\
    \hline

    \theutterance \stepcounter{utterance}  
    & & & \multicolumn{4}{p{0.6\linewidth}}{
        \cellcolor[rgb]{0.9,0.9,0.9}{
            \makecell[{{p{\linewidth}}}]{
                \texttt{\tiny{[P1$\langle$GM]}}
                \texttt{You are an expert in causal inference. The following question is not a typical commonsense query, but rather a meticulously designed question created by a professor specializing in causal inference, intended to assess the students' mastery of the course content.} \\
\texttt{Imagine a self{-}contained, hypothetical world with only the following conditions, and without any unmentioned factors or causal relationships: Yupt has a direct effect on muvq. Zupj has a direct effect on muvq. We know that yupt or zupj causes muvq. We observed an individual is not zupj. Would an individual is muvq if yupt instead of not yupt?} \\
\texttt{Start your answer with 'Yes' or 'No', followed by additional reasoning or evidence to support your explanation.} \\
            }
        }
    }
     \\ \\

    \theutterance \stepcounter{utterance}  
    & \multicolumn{4}{p{0.6\linewidth}}{
        \cellcolor[rgb]{0.561,0.737,0.561}{
            \makecell[{{p{\linewidth}}}]{
                \texttt{\tiny{[P1$\rangle$GM]}}
                \texttt{No} \\
\\ 
\texttt{**Reasoning:**} \\
\\ 
\texttt{To determine the answer, we must analyze the causal structure and the logical constraints provided in the hypothetical world.} \\
\\ 
\texttt{1.  **Causal Structure**:} \\
\texttt{    *   \$Yupt \textbackslash{}rightarrow muvq\$ (Direct effect)} \\
\texttt{    *   \$Zupj \textbackslash{}rightarrow muvq\$ (Direct effect)} \\
\texttt{    *   The statement "yupt or zupj causes muvq" implies that \$muvq\$ is the outcome of these causes. In causal inference contexts, unless specified as an exclusive OR (XOR) or a specific logical gate, "causes" usually implies a sufficient condition or a positive contribution. However, the critical piece of information is the observation.} \\
\\ 
\texttt{2.  **Observation**:} \\
\texttt{    *   We observed an individual is **not** \$Zupj\$ (\$\textbackslash{}neg Zupj\$).} \\
\\ 
\texttt{3.  **The Counterfactual Question**:} \\
\texttt{    *   "Would an individual be \$muvq\$ if \$Yupt\$ instead of not \$Yupt\$?"} \\
\texttt{    *   This asks for the counterfactual outcome \$Y\_\{yupt\}\$ given the background condition \$\textbackslash{}neg Zupj\$. Specifically, it asks if setting \$Yupt = 1\$ (true) results in \$muvq = 1\$ (true), compared to the baseline where \$Yupt = 0\$ (false).} \\
\\ 
\texttt{4.  **Analysis of} \\
            }
        }
    }
    & & \\ \\

    \theutterance \stepcounter{utterance}  
    & & & \multicolumn{2}{p{0.3\linewidth}}{
        \cellcolor[rgb]{0.9,0.9,0.9}{
            \makecell[{{p{\linewidth}}}]{
                \texttt{\tiny{[GM$|$GM]}}
                \texttt{No} \\
            }
        }
    }
    & & \\ \\

    \theutterance \stepcounter{utterance}  
    & & & \multicolumn{2}{p{0.3\linewidth}}{
        \cellcolor[rgb]{0.9,0.9,0.9}{
            \makecell[{{p{\linewidth}}}]{
                \texttt{\tiny{[GM$|$GM]}}
                \texttt{yes} \\
            }
        }
    }
    & & \\ \\

    \theutterance \stepcounter{utterance}  
    & & & \multicolumn{2}{p{0.3\linewidth}}{
        \cellcolor[rgb]{0.9,0.9,0.9}{
            \makecell[{{p{\linewidth}}}]{
                \texttt{\tiny{[GM$|$GM]}}
                \texttt{game\_result = LOSE} \\
            }
        }
    }
    & & \\ \\

\end{supertabular}
}

\end{document}
